\scanpage{55}
\mspage{42}
\margin{n°382}
\[
\begin{aligned}
s(tx)&=(st)x&&\text{si}\quad s,t\in G,\ x\in X\\
ex&=x&&\text{si}\quad x\in X\quad(e\text{ étant l'élément neutre de }G).
\end{aligned}
\]
On dit aussi que G \uline{opère à gauche} sur X ; on définit de même les lois d'opération à droite.

Soit $x\in X$, et soit $G_x$ son stabilisateur dans G. Supposons que l'application $g\longrightarrow g.x$ soit une \uline{subimmersion} (hypothèse automatiquement vérifiée si la caractéristique de K est 0). Alors $G_x$ est une sous-variété de groupe, et l'application de $G/G_x$ dans X obtenue par passage au quotient à partir de $g\longrightarrow g.x$ est une immersion. Si en outre l'orbite $Gx$ de x est localement fermée, et si la topologie de G admet une base dénombrable, $Gx$ est une sous-variété de X, et l'application $G/G_x\longrightarrow Gx$ est un isomorphisme.

4.12.6. Soit $\varpi:G\longrightarrow S$ une submersion, G et S étant deux variétés. Pour tout $s\in S$, soit $G_s=\varpi^{-1}(s)$. Supposons que chacun des $G_s$ soit muni d'une structure de groupe, vérifiant les deux conditions suivantes :

(i) L'application $e:S\longrightarrow G$ qui applique $s\in G$ en l'élément neutre $e(s)$ de $G_s$ est un morphisme.

(ii) L'application de $G\times_S G$ dans G définie par $(x,y)\longrightarrow x^{-1}y$ est un morphisme.

On dit alors que G est un \uline{groupe au-dessus de} S, ou encore un S-\uline{groupe}. Pour tout $s\in S$, $G_s$ est une sous-variété de G, et sa structure de variété est compatible avec sa structure de groupe.

Si G est un S-groupe, l'application $x\longrightarrow x^{-1}$ est un morphisme de G dans lui-même ; l'application $(x,y)\longrightarrow xy$ est un morphisme de $G\times_S G$ dans G.

4.12.7. Soit $\varpi:G\longrightarrow S$ un S-groupe, et soit $f:T\longrightarrow S$ un morphisme. Soit $G(T)$ l'ensemble des morphismes $\varphi:T\longrightarrow G$ tels que $\varpi\circ\varphi=f$. Si $\varphi,\psi\in G(T)$, soit $\varphi.\psi$ le composé de $(\varphi,\psi)_S:T\longrightarrow G\times_S G$ avec l'application $(x,y)\longrightarrow xy$ de $G\times_S G$ dans G ; on a $\varphi.\psi\in G(T)$ et la loi de composition ainsi définie fait de $G(T)$ un groupe d'élément neutre $e\circ f$.
\mspage{54}
